\section{Introduction}

Knowledge tracing (KT) estimates a learner's changing knowledge from a
chronological sequence of item responses \citep{corbett-anderson-1995,
abdelrahman-etal-2023}.  Its inputs normally associate each item with one or
more knowledge components (KCs), often through a binary Q-matrix.  That matrix
presupposes decisions that are easy to hide: what counts as a skill, how a
curriculum description becomes that representation, and whether an item really
exercises the selected components.  These are scientific choices, not file
formatting.

Grammar resources pose a sharp version of the problem.  The English Grammar
Profile (EGP) contains over 1,200 corpus-derived competence descriptions mapped
to CEFR levels \citep{okeeffe-mark-2017}.  A description can combine tense,
aspect, voice, polarity, clause type, and modality, state bounded alternatives,
or fall outside a selected ontology.  Treating every descriptor as one KC
obscures reusable operations; decomposing it without explicit rules can invent
unsupported distinctions or erase the source connection.

We introduce \system, a research harness that makes this path explicit.
Figure~\ref{fig:pipeline} separates source selection, model-assisted
normalization, canonicalization, realization, KC projection, item construction,
Q-matrix derivation, simulation, and KT evaluation.  Prompts, rulebooks,
policies, templates, and parameters remain replaceable versioned artifacts.

We ask:
\begin{enumerate}
  \item[RQ1] How completely and consistently can EGP prose be normalized into
        a closed grammar representation?
  \item[RQ2] How do factorized, interaction-augmented, and whole-cell policies
        change item--KC structure?
  \item[RQ3] Do all three resulting datasets support reproducible end-to-end KT
        checks, and what can those checks legitimately show?
\end{enumerate}

Our contributions are a conservative six-dimensional representation, a
provenance path from prose to Q-matrix entries, and an executable comparison of
three KC policies.  Unlike a human-data evaluation, this study cannot establish
that any KC is cognitively atomic or that predictive differences transfer to
learners; technical integrity and educational validity are kept separate.

\begin{figure*}[t]
  \centering
  \begin{tikzpicture}[
    >=Latex,
    node distance=1.8mm,
    box/.style={draw, rounded corners=1.5pt, align=center, minimum height=8mm,
      text width=1.62cm, font=\scriptsize, inner sep=2pt},
    model/.style={box, fill=orange!16, dashed},
    fixed/.style={box, fill=blue!10},
    truth/.style={box, fill=green!14},
    experiment/.style={box, fill=violet!11},
    arrow/.style={->, semithick},
    layer/.style={draw, rounded corners=2pt, inner sep=4pt}
  ]
    \node[fixed] (source) {grammar\\resource};
    \node[model, right=of source] (norm) {constrained\\normalisation};
    \node[fixed, right=of norm] (cells) {canonical\\GrammarCells};
    \node[truth, very thick, right=of cells] (kstar) {declared\\$\kstar$};
    \node[model, right=of kstar] (items) {generate, validate,\\curate items};
    \node[truth, right=of items] (qstar) {$\qstar$ + measurement\\audit};
    \node[truth, right=of qstar] (simulate) {baseline learner\\simulation};
    \node[fixed, very thick, right=of simulate] (dataset) {frozen dataset:\\items + responses};

    \draw[arrow] (source) -- (norm);
    \draw[arrow] (norm) -- (cells);
    \draw[arrow] (cells) -- (kstar);
    \draw[arrow] (kstar) -- (items);
    \draw[arrow] (items) -- (qstar);
    \draw[arrow] (qstar) -- (simulate);
    \draw[arrow] (simulate) -- (dataset);

    \node[experiment, below=9mm of dataset] (hide) {hide or perturb\\$\kstar,\qstar$};
    \node[experiment, left=of hide] (khat) {construct/select\\$\khat,\qhat$};
    \node[experiment, left=of khat] (evaluate) {prediction, state,\\structure};
    \node[experiment, left=of evaluate] (claims) {RQ2--RQ4 +\\robustness};

    \draw[arrow] (dataset) -- (hide);
    \draw[arrow] (hide) -- (khat);
    \draw[arrow] (khat) -- (evaluate);
    \draw[arrow] (evaluate) -- (claims);

    \node[layer, fit=(source)(dataset),
      label={[font=\scriptsize\bfseries]above:Layer A: frozen dataset construction}] {};
    \node[layer, fit=(claims)(hide),
      label={[font=\scriptsize\bfseries]below:Layer B: downstream hypotheses and tests}] {};
  \end{tikzpicture}
  \caption{\system{} separates dataset construction from downstream research.
  Generator truth is fixed before item responses exist.  Later discovery and
  misspecification experiments may hide or perturb it but never redefine the
  immutable baseline.  Model assistance is confined to the dashed linguistic
  normalisation and item-construction boxes.}
  \label{fig:pipeline}
\end{figure*}

